\documentclass[../../main.tex]{subfiles}

\begin{document}

\section{Number Theory}

There are quite a few important that I wish to discuss before
we do Diophantine equations. I don't think some of the contents
come out in the actual exam as many are focused on Diophantine
equations, but let's go over them anyways as a review.

\begin{definition}{}{}
A \textbf{residue class} of $a$ mod $n$ is the collection of
all numbers that are congruent to $a \pmod{n}$.
The \textbf{complete residue system} mod $n$ is the set of
exactly one representative from each of the $n$ residue classes.
The \textbf{reduced residue system} is a set of all elements
from the complete residue system that only contains the elements
that are relatively prime to $n$.
\end{definition}

An example for modulo $6$ is that the complete residue system
is $\{ 0, 1, 2, 3, \linebreak 4, 5 \}$ while the reduced residue
system is $\{ 1, 5 \}$. The reduced residue system for modulo $6$
does not have to be $\{ 1, 5 \}$, but $\{ 7, -1 \}$ too.

With this definition, we can also define a very special and
important function.

\begin{definition}{Euler's Totient Function}{Euler's Totient Function}
\textbf{Euler's Totient Function}, denoted as $\phi(n)$ is
the number of elements in the reduced residue system mod $n$.
\end{definition}

By definition, there are few ways to think about this function.

\begin{important}
$\phi(n)$ can be interpreted as the number of integers from
$1$ to $n$ that are relatively prime to $n$.
\end{important}

As you can see from our previous example of residue systems,
we can see that $\phi(6) = 2$. Another very important topic about
Euler's Totient function is its formula. Consider the following
equation for prime $p$.
\[
\phi\left( p^\alpha \right)
= p^\alpha - \left\lfloor \frac{p^\alpha}{p} \right\rfloor
= p^\alpha - p^{\alpha - 1}
= p^\alpha \left( 1 - \frac{1}{p} \right)
\]
This holds because the floor function, or Gauss function,
represents the number of elements that are multiples of $p$ and
$p$ is relatively prime to other numbers that are not multiple
of $p$. Before deriving the formula, let's briefly take a look
at its multiplicative property.

\begin{important}
$\phi(ab) = \phi(a) \phi(b)$ for $(a, b) = 1$.
\end{important}

There are few ways to prove this property, but let's prove with
the following arrangement.
\[
\begin{array}{ccccc}
    1      & 2      & 3      & \cdots & b  \\
    b + 1  & b + 2  & b + 3  & \cdots & 2b \\
    2b + 1 & 2b + 2 & 2b + 3 & \cdots & 3b \\
    \vdots & \vdots & \vdots & \ddots & \vdots \\
    \cdots & \cdots & \cdots & \cdots & ab
\end{array}
\]
As you can see, $\phi(b)$ is the number of integers that are
relatively prime to $b$ in the first row. In other words,
there exists exactly $\phi(b)$ columns that contain the integers
that are relatively prime to $b$. Let $j$ be a relatively prime
integer to $b$. Because $(a, b) = 1$, it is self-evident that
there exists $\phi(a)$ numbers that are relatively prime to $a$
in $j^\text{th}$ column. Because there exists $\phi(b)$ such columns,
the total number of integers that are relatively prime to both
$a$ and $b$, or relatively prime to $ab$, is $\phi(a) \phi(b)$ and
$\phi(ab) = \phi(a) \phi(b)$ holds for $(a, b) = 1$.

\begin{important}
Continuing with this property, we can see that the same applies
if we prime factorize $n$ as $n = p_1^{\alpha_1} p_2^{\alpha_2}
\cdots p_k^{\alpha_k}$ for prime $k$. In other words, we can write
the formula as the following.
\[
\phi(n)
= p_1^{\alpha_1} \left( 1 - \frac{1}{p_1} \right)
    \cdots
    p_k^{\alpha_k} \left( 1 - \frac{1}{p_k} \right)
= n \prod_{i=1}^k \left( 1 - \frac{1}{p_i} \right)
\]
\end{important}

There is also very important property
\[
\sum_{d \mid n} \phi(d) = n
\]
but maybe we can discuss this later.

Continuing with the totient function, we can derive two important
theorems.

\begin{theorem}{Euler's Theorem}{Euler's Theorem}
$a^{\phi(n)} \equiv 1 \pmod{n}$ if $(a, n) = 1$.
\end{theorem}
\begin{proof}
Let $\{ r_1, r_2, \ldots, r_{\phi(n)} \}$ be the set of integers
that are relatively prime to $n$. Note that for positive integer $a$
such that $(a, n) = 1$, $\{ ar_1, ar_2, \ldots, \linebreak
ar_{\phi(n)} \}$ is the same set since $ar_i \not\equiv ar_j \pmod{n}$
for $i \neq j$.

Therefore,
\[
r_1 r_2 \ldots r_{\phi(n)}
\equiv ar_1 ar_2 \ldots ar_{\phi(n)} \pmod{n}
\]
holds and $a^{\phi(n)} \equiv 1 \pmod{n}$.
\end{proof}

The next theorem Fermat's Little Theorem can be proven with multiple
ways, but let's use Euler's Theorem to prove it. They are frequently
introduced together and some say Fermat's Little Theorem is a
corollary of Euler's Theorem, but let's consider it as a standalone
theorem because there are plenty of other ways to prove the theorem.

\begin{theorem}{Fermat's Little Theorem}{Fermat's Little Theorem}
For prime $p$ and positive integer $a$ such that $(a, p) = 1$,
$a^{p-1} \equiv 1 \pmod{p}$.
\end{theorem}
\begin{proof}
By Euler's Theorem, $a^{\phi(p)} \equiv 1 \pmod{p}$. Because
$\phi(p) = p - 1$, the theorem holds.
\end{proof}

Another very common way of representing this theorem is
$a^p \equiv a \pmod{p}$. Before we move on to the final theorem
that we will discuss in this section, let's take a look at an
example of how they are used.

\begin{problem}{}{}
Show that $8^{43} \equiv 2 \pmod{10}$.
\end{problem}
\begin{solution}
We can use CRT for this, but let's use Fermat's Little Theorem here.
How can we use it if $(8, 10) \neq 1$? Well, we can force it.
Consider the following congruences with $8^{43} \equiv 0 \pmod{2}$
in mind.
\begin{align*}
    8^4 &\equiv 1 \pmod{5} \\
    8^{43} &\equiv 8^3 \equiv 2^9 \equiv 2 \pmod{5}
\end{align*}
In other words, $8^{43} = 5m + 2 = 2n$ for some integers $m$ and $n$.
Because $2 \mid m$, $8^{43} \equiv 2 \pmod{10}$.
\end{solution}

Like this, we can split the modulo $n$ for our appropriate need.
Before moving on to the final one and concepts related to the
theorem, here is another example.

\begin{problem}{}{}
Show that $2009^{2008^{2007}} \equiv 81 \pmod{1000}$ \cite{MATH01-5}.
\end{problem}
\begin{solution}
First, notice that $(2009, 1000) = 1$ and $\phi(1000) = 400$.
Therefore, $2009^{400} \equiv 1 \pmod{1000}$ and the following
congruence holds for some $n$ and $k$.
\[
2009^{2008^{2007}}
\equiv 2009^{400n + k}
\equiv 2009^{k}
\equiv 9^k \pmod{1000}
\]
Continuing, note that $2008^{2007} \equiv k \pmod{400}$.


CONTINUE
\end{solution}




\end{document}


